% ECONOMICS_PROBLEM: {"id": "Q58-273", "title": "Political durability and design of carbon pricing", "jel_code": "Q58", "source_id": "OP-0116"}
% !TeX program = xelatex
\documentclass[11pt,a4paper]{article}
\usepackage[margin=23mm]{geometry}
\usepackage{fontspec}
\setmainfont{Latin Modern Roman}
\usepackage{amsmath,amssymb,mathtools}
\usepackage{xcolor,enumitem,booktabs,longtable,array,xurl,hyperref}
\definecolor{muted}{HTML}{53636C}
\hypersetup{colorlinks=true,urlcolor=blue,linkcolor=blue}
\setlength{\parindent}{0pt}
\setlength{\parskip}{5pt}
\newcommand{\R}{\mathbb R}
\newcommand{\E}{\mathbb E}
\newcommand{\N}{\mathbb N}
\newcommand{\ind}{\mathbf 1}
\newcommand{\Var}{\operatorname{Var}}
\newcommand{\Cov}{\operatorname{Cov}}
\newcommand{\supp}{\operatorname{supp}}
\DeclareMathOperator*{\argmin}{arg\,min}
\DeclareMathOperator*{\argmax}{arg\,max}
\newcommand{\entrymeta}[3]{{\sffamily\small\color{muted}\textbf{\texttt{#1}}\quad|\quad #2\par #3\par}}
\newcommand{\Problemref}[1]{\texttt{#1}}
\newcommand{\JELref}[2]{\texttt{#2} (JEL \texttt{#1})}
\begin{document}
\section*{Political durability and design of carbon pricing}
\textbf{Problem ID:} \texttt{Q58-273}\quad \textbf{JEL:} \texttt{Q58}\par
% BEGIN ECONOMICS STATEMENT
\label{problem:OP-0116}
\label{atlas:prob:C6}

\noindent\textbf{Original identifier:} C6 (atlas item 116). \textbf{Field:} Environmental and climate economics. \textbf{Classification:} editorial research family with established restricted solutions; current open status not certified.

\subsection*{Definitions and assumptions}

At dates $t=0,\ldots,H$, a carbon-pricing regime $R$ specifies tax $\tau_t$ in real currency per tonne of $\mathrm{CO}_2$, coverage, exemptions, border rules and household rebates $b_{h,t}$. Define taxable emissions $E_t$ separately from total emissions. A government budget includes revenue $\tau_tE_t$, rebate payments, administration costs and other financing sources. A political-economic state $x_t$ records household assets, technologies, preferences and incumbent policy. Under a stated voting/lobbying institution, private behavior and policy continuation form an equilibrium $e\in\mathcal E(R,P)$ for law $P$. Repeal time $T_{\rm rep}$ is the first date the declared core regime ceases.

\subsection*{Formal research question}

Determine policy designs that jointly satisfy fiscal feasibility, emissions objectives and political durability. With supplied social weights $\omega_h$, utilities $u_h$, discount factors $\beta_t$ and consumption $c_{h,t}(e)$, define $W(R,e,P)=\mathbb E_P[\sum_t\beta_t\sum_h\omega_hu_h(c_{h,t}(e))]$, including monetized or explicitly modeled climate damages consistently. For an ambiguity class $\mathcal P$ and durability target $\kappa\in(0,1)$, seek a rule maximizing worst-case welfare subject to
\[
\inf_{P\in\mathcal P,e\in\mathcal E(R,P)}
\Pr_P(T_{\rm rep}>H)\ge\kappa
\]
and a stated cumulative-emissions bound. Characterize how rebate visibility, commitment limits, exemption-induced rents and border responses change the feasible welfare-durability frontier. Specify equilibrium selection if performance is evaluated at one equilibrium instead of over all equilibria.

\subsection*{Interpretation and scope}

Durability is not a synonym for popularity in a survey; actual support depends on realized prices, transfers, trust and organized opposition. A durable zero tax need not meet the emissions constraint. Exemptions affect marginal incentives and general-equilibrium incidence, while lump-sum rebates can alter political support without mechanically lowering the marginal tax rate. Dynamic policy credibility changes investment before elections, so a one-period static willingness-to-pay estimate is not enough to identify continuation probabilities. The welfare objective embeds intergenerational discounting and regional weights rather than presuming a universally agreed efficient path. Constitutional or legal commitment is modeled as an institutional restriction, not assumed enforceable by notation.

\subsection*{Source audit and status}

[S1] is a verified completion using Metcalf's 2019 Brookings article; his 2019 book is another possible intended source. The Greenstone (2019) citation remains unresolved [S2] after inspecting the author's research bibliography. [S3] develops political economy of green transitions. The broad OECD taxation link supplies institutional background, not an exact theoretical source. The durability frontier is editorial and is not certified by these citations as a published or currently open conjecture.

\subsection*{References}

\begin{enumerate}[label={[\textnormal{S}\arabic*]},leftmargin=*]

\item Gilbert E. Metcalf (2019). \emph{On the Economics of a Carbon Tax for the United States}. Brookings Papers on Economic Activity, Spring, 405–484. \url{https://www.brookings.edu/articles/on-the-economics-of-a-carbon-tax-for-the-united-states/}. \textbf{Locator:} Publisher citation, lines 163–168. \textbf{Support and limits:} Relevant verified completion; author has another 2019 carbon-tax book, so exact intended source is not certain.

\item \textbf{Unresolved original citation:} Greenstone (2019). Author research list searched; no unique intended 2019 carbon-pricing publication identified from author-year alone. Candidate/verification route: \url{https://michaelgreenstone.com/research/}.

\item Timothy Besley and Torsten Persson (2023). \emph{The Political Economics of Green Transitions}. Quarterly Journal of Economics 138(3), 1863–1906. \url{https://academic.oup.com/qje/article/138/3/1863/7000849}. \textbf{Locator:} Publisher or institutional abstract and bibliographic record. \textbf{Support and limits:} Dynamic political economy with endogenous values and technologies; does not imply every coalition sustains optimal pricing.

\end{enumerate}

\subsection*{Common conventions: Source mathematical conventions}

Unless an entry states otherwise, agent and firm sets are finite. State and action spaces are measurable, strategies and outcome maps are measurable, probability laws are specified, and expectations exist for the targets under discussion. Infinite-horizon discount factors lie in $(0,1)$ and discounted objectives are finite. Norms, tolerances and complexity input sizes must be specified for computational comparisons. Constants or primitives that appear locally are inputs, not empirical facts asserted by this catalogue.

\textbf{Identification.} A statistical model $m\in\mathcal M$ implies an observable law $P_m$ and target $\theta(m)$. At law $P$, its identified set is
\[
\Theta(P;\mathcal M)=\{\theta(m):m\in\mathcal M,\ P_m=P\}.
\]
Point identification holds when this set is a singleton, or equivalently when $P_{m_1}=P_{m_2}$ implies $\theta(m_1)=\theta(m_2)$ for all compatible models. Sharp bounds mean bounds attained, or approached when the boundary is not attained, by compatible models. Writing this set is a definition, not a solution: the substantive task is to establish informative restrictions and derive or estimate the set. An unrestricted-confounding model may give only trivial bounds.

\textbf{Inference.} Identification, estimability and valid uncertainty quantification are separate questions. Fix $\alpha\in(0,1)$. A confidence procedure $C_n$ over a declared law class $\mathcal P$ has asymptotic uniform coverage at level $1-\alpha$ when
\[
\liminf_{n\to\infty}\inf_{P\in\mathcal P}\Pr_P\{\theta(P)\in C_n\}\geq1-\alpha.
\]
For set-identified targets the coverage event must be defined separately, for example coverage of the full identified set. If an entry uses identification-indexed models instead of uniquely indexed laws, the same coverage convention is applied over those models. Pointwise limits do not establish uniform validity, and asymptotic validity does not guarantee finite-sample precision. Sample sizes, dependence structures and weak-identification sequences are part of each application.

\textbf{Counterfactuals.} $Y_i(a)$ is a potential outcome under a specified intervention, including its timing, implementation and versions. It is not $Y_i$ conditional on voluntarily choosing $a$. A full intervention vector permits interference; shorthand scalar treatment notation does not silently impose no interference. Assignment, selection, attrition, anticipation and measurement must be specified. A claim about an instrument's local effect is not automatically a population policy effect.

\textbf{Economic equilibrium.} A counterfactual model includes best responses, constraints, feasibility, market clearing, state transitions and expectations consistent with the stated information. When several equilibria exist, either a declared selector is an input or the target is set-valued. Comparisons across policy regimes must use the stated selection convention. Reduced-form relationships are not presumed invariant after an equilibrium-changing intervention.

\textbf{Optimization and welfare.} Optimization is over the stated feasible set. If attainment is not established, $\sup$ or $\inf$ and $\varepsilon$-optimal choices replace an asserted maximizer or minimizer. For example, with value $V=\sup_{a\in\mathcal A}W(a)<\infty$, an $\varepsilon$-optimal set is $\{a\in\mathcal A:W(a)\geq V-\varepsilon\}$ for $\varepsilon>0$. Benefits, costs, risk and health or environmental outcomes can be aggregated only using declared units and conversion weights. Interpersonal utility weights, the treatment of fiscal transfers and foreign profits, discounting, and the behavioral welfare criterion are explicit inputs. A comparative-static derivative additionally requires existence and appropriate differentiability of the selected counterfactual.

\textbf{Sources.} Local labels [S1], [S2], and so forth restart in every question. Locators identify the relevant abstract, model, section or result at the level actually checked; a bibliographic or abstract-level verification is not represented as inspection of an entire proof. Working papers are distinguished from later publications and changing versions. Source summaries are paraphrased. All URL checks and editorial formulations refer to the audit date stated above.
% END ECONOMICS STATEMENT
\end{document}
